\begin{table}[H]
\centering
\caption{Sensitivity to the 2024 adoption cohort and to the 2024 CRS vintage year}
\label{tab:att_drop2024}
% latex table generated in R 4.6.0 by xtable 1.8-8 package
% Thu Sep 17 13:21:21 2026
\begingroup\small
\adjustbox{max width=\textwidth}{%
\begin{tabular}{llllllll}
  \toprule
Specification & ATT & SE & $t$ & Pre-trend $p$ & $N$ & Recipients & Treated \\ 
  \midrule
Baseline (headline specification) & 0.2928** & (0.1241) & 2.360 & 0.737 & 2,030 & 127 & 40 \\ 
  Drop the 2024 adoption cohort & 0.2492** & (0.1256) & 1.984 & 0.480 & 1,902 & 119 & 32 \\ 
  Drop calendar year 2024 (and the 2024 cohort) & 0.3322*** & (0.1195) & 2.781 & 0.480 & 1,783 & 119 & 32 \\ 
   \bottomrule
\end{tabular}
}
\endgroup
\par\vspace{4pt}
\begin{minipage}{\linewidth}\footnotesize
Notes: Outcome: log(adaptation commitments). Row 2 drops the 2024 adoption cohort; row 3 also drops calendar year 2024 for every recipient, which removes that cohort's post-treatment period entirely. Baseline row read from the stored headline fit. Unless noted: headline specification (CS\,(2021) DR, never-treated controls, WGI GE + log population, cohorts $\geq$ 5), multiplier-bootstrap SE (999 reps, seed 1242). Pre-trend $\chi^2$ / $p$: joint Wald test on the aggregated pre-treatment event-time coefficients ($-5 \leq e \leq -2$; 4 restrictions), using the influence-function covariance of the dynamic aggregation from the analytical (non-bootstrap) fit; a generalized inverse is used if the block is singular (analytical twin fit).\par
Source: OECD CRS (Rio adaptation markers); UNFCCC NAP Central.
\end{minipage}
\end{table}
